\chapter{Conclusioni}

Questa tesi ha avuto come obiettivo la costruzione di un modello matematico in grado di descrivere l'evoluzione, seduta dopo seduta, di un gruppo di judoka nel corso di un periodo di allenamento, allo scopo di analizzare gli infortuni e l'incidenza della prevenzione. \par

Gli atleti sono suddivisi dal modello in quattro popolazioni: i sani, coloro che hanno subito infortuni pregressi, gli infortunati gravi, impossibilitati ad allenarsi, e i malconci, che continuano ad allenarsi nonostante un infortunio lieve o una guarigione incompleta. A queste popolazioni si affiancano due variabili che rappresentano la fatica media del gruppo, nella componente fisica e in quella psicologica. Cinque funzioni governano il passaggio da una popolazione all'altra: esse descrivono l'arrivo di nuovi atleti, gli abbandoni, gli infortuni, il carico di allenamento e il recupero. Tra queste, la funzione di infortunio riveste il ruolo principale, poiché mette in relazione il rischio con la fatica del gruppo e con la densità di atleti presenti sul tatami. \par

Per ogni funzione abbiamo adottato una forma analitica semplice, dotata di parametri interpretabili, e ne abbiamo esaminato le proprietà. In mancanza di dati reali, abbiamo proceduto a una calibrazione per plausibilità: dopo aver stabilito alcuni requisiti che la dinamica del gruppo deve soddisfare, abbiamo individuato, tramite la libreria SciPy, i parametri più prossimi a una stima iniziale che li rispettano. Ne è derivato uno scenario di riferimento che rappresenta un periodo di allenamento intenso privo di prevenzione. \par

Abbiamo infine inserito nel modello un controllo, che esprime l'efficacia della prevenzione attuata dall'allenatore e riduce il rischio di infortunio per tutte le popolazioni. Abbiamo quindi simulato una strategia di prevenzione commisurata al carico, più incisiva nei giorni di randori, ridotta della metà nei giorni di tecnica e del tutto assente nelle sedute leggere, confrontandola con lo scenario privo di prevenzione.

\section{Risultati principali}

Nello scenario senza prevenzione, nell'arco di cinquanta sedute si registrano circa $428$ infortuni e il gruppo perde il $15\%$ dei propri atleti. La quasi totalità degli abbandoni, circa l'$87\%$, riguarda atleti infortunati o malconci, e al termine del periodo circa un terzo degli atleti rimasti non si trova in condizioni ottimali. \par

Le simulazioni hanno messo in luce tre aspetti fondamentali.
\begin{itemize}
    \item \textbf{La fatica accresce il rischio.} Il rischio di infortunio aumenta con la fatica accumulata, e difatti gli infortuni risultano più frequenti nella seconda metà del periodo. Il nesso tra carico di allenamento, fatica e rischio di infortunio costituisce un tema centrale della medicina dello sport~\cite{soligard2016}. La scelta della funzione di recupero si è dimostrata determinante: qualora il recupero cali troppo rapidamente, la fatica cresce senza limite, il rischio raggiunge il valore massimo e il carico di allenamento perde ogni effetto sugli infortuni.
    \item \textbf{I malconci costituiscono la popolazione più critica.} Benché siano meno numerosi dei sani, essi producono circa il $41\%$ degli infortuni, in quanto popolazione più esposta al rischio. Essi alimentano inoltre un ciclo nel quale gli infortuni lievi conducono a nuovi infortuni gravi.
    \item \textbf{Gli infortuni determinano gli abbandoni.} Gli infortunati e i malconci abbandonano con molta maggiore facilità rispetto agli atleti sani. Ridurre gli infortuni equivale dunque anche a mantenere unito il gruppo.
\end{itemize}

La prevenzione proporzionale al carico, con una prevenzione media di circa $0.3$ per seduta, determina una riduzione degli infortuni del $31\%$, da circa $428$ a circa $296$, e degli abbandoni del $18\%$, da circa $124$ a circa $102$. L'effetto risulta particolarmente marcato sui malconci, i cui infortuni calano quasi della metà, poiché la prevenzione impedisce che molti atleti diventino malconci. Infine, l'effetto della prevenzione cresce in modo regolare con l'intensità: anche una prevenzione leggera produce un beneficio apprezzabile. La \autoref{fig:riepilogo} riassume il confronto tra i due scenari.

\begin{figure}[ht]
    \centering
    \includegraphics[width = \textwidth]{conclusioni/riepilogo.pdf}
    \caption{Infortuni e abbandoni totali, con e senza prevenzione}\label{fig:riepilogo}
\end{figure}


\section{Limiti del modello}

Il modello costituisce una semplificazione della realtà, e talune scelte compiute per mantenerlo semplice ne circoscrivono la validità. È opportuno esplicitarle, sia per interpretare correttamente i risultati sia per individuare possibili miglioramenti.

\paragraph{Assenza di dati reali}

Il limite principale risiede nel fatto che i parametri non sono stati stimati a partire da dati reali. La calibrazione per plausibilità assicura che la dinamica del gruppo sia ragionevole, ma il risultato dipende da due scelte soggettive: i valori iniziali dei parametri e le soglie dei requisiti, quali il calo della popolazione o il picco di infortunati gravi. Scelte differenti condurrebbero a uno scenario diverso. \par

Per tale ragione i numeri ottenuti, come i circa $428$ infortuni o la riduzione del $31\%$ attribuibile alla prevenzione, non devono essere interpretati come previsioni di quanto accade in una palestra reale. Essi assumono valore soprattutto nel confronto tra scenari differenti: risulta ad esempio più affidabile l'affermazione secondo cui la prevenzione riduce gli infortuni dei malconci più di quelli dei sani, rispetto al valore esatto della riduzione. Inoltre, combinazioni diverse di parametri possono generare traiettorie simili, e in assenza di dati non è possibile stabilire quale sia quella corretta. Gli studi sugli infortuni nel judo reperibili in letteratura descrivono soprattutto la frequenza e la tipologia degli infortuni \cite{pocecco2013}, ma non seguono l'evoluzione di un gruppo di atleti seduta dopo seduta, che è quanto servirebbe per stimare i parametri del modello.

\paragraph{La fatica è identica per tutti}

Il modello impiega un'unica fatica media per l'intero gruppo, scelta giustificata dalle ipotesi di omogeneità, mescolamento e simmetria. Nella realtà, tuttavia, la fatica varia da atleta ad atleta: un malconcio è verosimilmente più affaticato di un sano, e un atleta appena rientrato da un infortunio lo è meno di uno che si è allenato per l'intero periodo. Il modello non è in grado di rappresentare tali differenze, né la possibilità che un allenatore assegni carichi differenti ad atleti differenti. Inoltre, nelle simulazioni il carico fisico e quello psicologico coincidono, per cui la distinzione tra le due componenti della fatica, pur presente nel modello, non è stata sfruttata.

\paragraph{Il recupero e l'orizzonte temporale}

La funzione di recupero adottata tende a zero con il trascorrere del tempo. Su un orizzonte di cinquanta sedute ciò non rappresenta un problema, poiché il recupero diminuisce lentamente e al termine del periodo vale ancora circa $0.15$. Su orizzonti più estesi, invece, la fatica crescerebbe senza limite e il rischio raggiungerebbe il valore massimo, con gli stessi problemi osservati per il recupero che decade rapidamente. Per orizzonti lunghi sarebbe pertanto necessario introdurre un valore minimo del recupero.

\paragraph{Ingressi e abbandoni}

Ingressi e abbandoni dipendono unicamente dal tempo, e la popolazione a cui si applicano interviene soltanto attraverso un moltiplicatore. Nella realtà, al contrario, la decisione di abbandonare può dipendere dalla fatica, dalla durata di un infortunio o dal clima del gruppo, mentre gli ingressi possono dipendere dalla dimensione del gruppo o dal periodo dell'anno. Il modello coglie tali differenze solo indirettamente, mediante i diversi moltiplicatori di abbandono delle quattro popolazioni.

\paragraph{Il controllo}

Il controllo è stato modellato nella forma più semplice possibile: riduce il rischio di infortunio di una frazione determinata, identica per tutte le popolazioni. Tale scelta trascura diversi aspetti:
\begin{itemize}
    \item la prevenzione non comporta alcun costo per l'allenamento: il tempo dedicato al riscaldamento o al rinforzo non riduce il carico né il beneficio della seduta;
    \item la prevenzione non modifica la gravità degli infortuni né la velocità di guarigione, sebbene nella realtà alcune misure, come la fisioterapia, potrebbero incidere su di esse;
    \item la prevenzione è identica per tutte le popolazioni, mentre un allenatore potrebbe riservare maggiore attenzione proprio ai malconci, che risultano i più a rischio;
    \item la strategia di prevenzione è stata scelta a priori e non ottimizzata: non è noto se esista una strategia migliore a parità di prevenzione complessiva.
\end{itemize}

\paragraph{Un modello deterministico}

Il modello è infine deterministico e tratta le popolazioni come numeri reali. Gli infortuni, tuttavia, sono eventi casuali: due periodi di allenamento identici possono presentare un numero di infortuni molto differente. Il modello descrive soltanto l'andamento medio, ed è tanto più affidabile quanto più numeroso è il gruppo. Per gruppi piccoli, come quelli di molte palestre, le fluttuazioni casuali possono assumere rilevanza pari a quella dell'andamento medio.

\section{Sviluppi futuri}

I limiti discussi indicano alcune direzioni naturali per lo sviluppo di questo lavoro.
\begin{itemize}
    \item \textbf{Raccogliere dati reali.} Anche pochi dati, quali il numero di iscritti, di abbandoni e di infortuni in una stagione in una singola palestra, consentirebbero di stimare i parametri anziché sceglierli, e di verificare se il modello descriva correttamente un gruppo reale.
    \item \textbf{Ricercare la strategia ottima.} Anziché scegliere la strategia a priori, si potrebbe individuare quella che riduce maggiormente gli infortuni a parità di prevenzione complessiva, oppure quella che raggiunge un determinato obiettivo con il minimo sforzo, avvalendosi ad esempio delle tecniche del controllo ottimo o dell'apprendimento per rinforzo.
    \item \textbf{Prevenzione mirata.} Consentire una prevenzione differenziata per ciascuna popolazione permetterebbe di studiare strategie rivolte in particolare ai malconci, che i risultati indicano come la popolazione più critica.
    \item \textbf{Il carico come controllo.} Oltre alla prevenzione, anche il carico di allenamento costituisce una scelta dell'allenatore. Considerarlo come seconda variabile di controllo permetterebbe di ricercare un equilibrio tra un allenamento efficace e un numero contenuto di infortuni.
    \item \textbf{Un modello stocastico.} L'introduzione della casualità degli infortuni permetterebbe di descrivere anche gruppi piccoli e di stimare non solo il numero medio di infortuni, ma anche la sua variabilità da un periodo all'altro.
\end{itemize}

In conclusione, il modello proposto è semplice, ma consente già di comprendere alcuni meccanismi rilevanti, quali il legame tra fatica e rischio, il ruolo dei malconci e il nesso tra infortuni e abbandoni, e di confrontare quantitativamente strategie di prevenzione differenti. Con l'aggiunta di dati reali e di alcune delle estensioni proposte, esso potrebbe divenire uno strumento utile per supportare gli allenatori nella progettazione di allenamenti più sicuri.